\section{Conclusiones y trabajo futuro}\label{sec:conclusiones}

\kyrian{} demuestra que un planificador de viajes con IA generativa puede ser
\textbf{conversacional y verificable a la vez}. Los cinco objetivos se cumplen con evidencia
medible:

\begin{itemize}
  \item \textbf{O1, fundamentación}: ningún lugar llega al itinerario sin un documento del corpus
        detrás. Está garantizado por construcción (herramientas que solo aceptan ids conocidos y
        tarjetas reconstruidas desde el corpus) y verificado en las 324 comprobaciones de seguridad
        de la evaluación del agente; las respuestas fundamentadas en el corpus obtienen además un 4,72/5 de un juez validado frente a un conjunto de referencia.
  \item \textbf{O2, edición conversacional}: el agente con herramientas duplicó los escenarios
        superados frente al planificador por intenciones (del 44,8 al 91,7\,\%) y, con el modelo de
        producción, supera el 92,6\,\%; en producción, la capa de identificadores conocidos
        descartó 26 ids elegidos fuera de los candidatos sin que ninguno llegara al viajero.
  \item \textbf{O3, corpus abierto y medido}: seis ciudades y 33.264 documentos de texto con
        licencia comprobada, puerta de calidad de 11 umbrales, R@10 de 0,768 a solo 0,015 del techo exacto,
        y un procedimiento agéntico (\code{/add-city}) para añadir ciudades.
  \item \textbf{O4, siempre en línea al menor coste}: arquitectura 100\,\% \emph{serverless} en
        AWS por unos 5\,€/mes estimados en reposo y alrededor de un céntimo por turno de
        conversación, con una cuota diaria por usuario; en 30 días de producción, ningún error de
        servidor.
  \item \textbf{O5, ingeniería medible}: IaC, CI con 13 trabajos y más de 2.000 pruebas automáticas,
        contenedores, tres niveles de observabilidad y un proceso de entrega con agentes
        documentado y automatizado.
\end{itemize}

Más allá del producto, el proyecto aporta un método: \textbf{las decisiones se apoyan en
mediciones} (el experimento que validó la base vectorial, el cambio de modo del índice medido antes
y después, la regla de no regresión del corpus y la evaluación del agente como requisito de cada
cambio del planificador) y \textbf{los agentes de IA trabajan dentro de límites explícitos}: \emph{briefs} escritos,
puertas automáticas y decisiones irreversibles reservadas a personas.

\subsection{Oportunidades de mejora}

\begin{enumerate}[label=\textbf{M\arabic*}]
  \item \textbf{Recuperación híbrida y reescritura de consultas}: las intenciones sin palabras
        clave («niños y lluvia») son el principal margen de mejora; combinar búsqueda léxica y
        vectorial, o reformular la consulta con el modelo, es el siguiente paso, ya previsto en la
        arquitectura objetivo (OpenSearch Serverless, Fig.~\ref{fig:objetivo}) uno de los dos criterios fijados para revisar la base vectorial (§\ref{sec:spike}).
  \item \textbf{Evaluación continua y más amplia}: llevar la evaluación del agente y la del juez a
        la CI con un presupuesto acotado; extender el juez de fundamentación a las respuestas del
        propio agente y repetir la evaluación sobre la pila completa (base vectorial y
        Bedrock Guardrails); y mejorar el juez, medido sobre el conjunto de referencia ya construido (pidiéndole que
        enumere cada dato concreto antes de puntuar, o con un modelo más capaz), con revisión humana
        de sus etiquetas.
  \item \textbf{Evaluación con usuarios}: un estudio con usuarios reales (tasa de éxito, tiempo y
        cuestionario de usabilidad SUS, \emph{System Usability Scale}) completaría las métricas automáticas.
  \item \textbf{Frescura del corpus}: reconstrucción programada y detección de cambios en las
        fuentes, con la misma puerta (Fig.~\ref{fig:objetivo}).
  \item \textbf{Más ciudades e idiomas}: con \code{/add-city}, un destino nuevo cuesta horas de
        agente y céntimos de cómputo.
  \item \textbf{Operación de producción}: evolución a la arquitectura objetivo (§\ref{sec:prod}),
        con EKS en subredes privadas, Aurora PostgreSQL, agente en Bedrock AgentCore, tres
        entornos y alarmas.
  \item \textbf{Monetización}: enlaces de afiliación en las tarjetas, fotos con licencia para la
        oferta B2B y exposición del planificador como aplicación dentro de asistentes generalistas
        (MCP), un canal que ya usan las grandes agencias en línea~\cite{lee2025chatgptapps}.
\end{enumerate}

Ninguna de estas mejoras exige cambiar el diseño: todas se apoyan en los puertos, el corpus
versionado y las evaluaciones que ya existen. Esa es la principal conclusión del proyecto: un
planificador con IA en el que \textbf{el modelo decide y el corpus describe} no solo es más
fiable; también es más barato de operar, más fácil de medir y más sencillo de hacer crecer.
